\section{Compare input bytes, output stores and simulated work}
Use several image families: the two-color sprite, uniform color, small component changes,
a sequence with frequent cache collisions, and unrelated RGBA values. Record the pixel
count, encoded byte count and executed simulator instructions for a named function
boundary. A small file does not prove a small instruction count, and a short input
stream does not imply few output stores when its runs expand many pixels.

Compare the literal baseline with the compressed policy on the same image. Count all
header work, cache initialization, capacity checks, calls and return paths. For decode,
include the validation pass when timing the complete interface. Measuring only the
second reconstruction pass answers a different question. Report that separate boundary
explicitly if it helps explain where the work occurs.

Keep simulator accuracy fixed while optimizing. A wider load must not read beyond the
selected input merely because the mapped arena has spare bytes. A cached state must
not survive a new invocation. A fast run expander must not publish pixels before
validation or overrun the image's logical extent. The private tests and public state
contracts must continue to cover those boundaries after an optimization.

The native cache is only 256 bytes. Before adding a more elaborate representation,
measure clearing it, updating it and reconstructing the stream separately. The full
program's file operations and IDE display work have different costs from the core
instruction count. None of those measurements should be presented as an IDE launch
latency without measuring that path.
